% GENERATED FILE. Edit canonical lessons, then run npm run generate:books.
% canonical-source-hash: fnv1a64:46d56b658da1e8b5
% canonical-lessons: MW-C158-polar, MW-C158-answer, MW-C158-kad, MW-C158-olakhno, MW-W48-o, MW-R158-yes-or-no

\chapter{Yes or No?}
\label{ch:mw-a2-yes-or-no}
\hlchaptermodality{158}

\begin{chapteropening}
\textbf{By the end of this chapter:} \emph{I can ask a yes-or-no question, answer it politely, ask when, and say that I recognise something.}

Complete MW-R158-yes-or-no: ask and answer yes-or-no questions politely, ask when, and recognise a song.
\end{chapteropening}

\section[\texorpdfstring{\mw{आप} \mw{चाय} \mw{पीवो} \mw{हो}? (āp chāy pīvo ho?)}{आप चाय पीवो हो? (āp chāy pīvo ho?)}]{\mw{आप} \mw{चाय} \mw{पीवो} \mw{हो}? — a yes-or-no question}

\label{lesson:MW-C158-polar}



\begin{quote}
Say \emph{you drink tea}, to \textbf{\mw{आप}}: \textbf{\mw{आप} \mw{चाय} \mw{पीवो} \mw{हो।}}
\end{quote}

\begin{grammarlens}[title={Grammar Lens: the voice does it}]
Say the same sentence again and let your voice rise at the end. That is the whole question:

\begin{quote}
\textbf{\mw{आप} \mw{चाय} \mw{पीवो} \mw{हो}?}
\end{quote}

\begin{quote}
\emph{āp chāy pīvo ho?} — Do you drink tea?
\end{quote}

\begin{quote}
\textbf{\mw{वो} \mw{बीमार} \mw{है}?}
\end{quote}

\begin{quote}
\emph{vo bīmār hai?} — Is he ill?
\end{quote}

Nothing moves and nothing is added. In print, only the end mark changes: \textbf{?} for \textbf{\mw{।}}.
\end{grammarlens}

\subsection*{Your turn}
Say these aloud:

\begin{itemize}
\raggedright
  \item \textbf{\mw{आप} \mw{दूध} \mw{पीवो} \mw{हो}?} — \emph{āp dūdh pīvo ho?} — Do you drink milk?
  \item \textbf{\mw{राम} \mw{अठै} \mw{है}?} — \emph{rām aṭhai hai?} — Is Ram here?
  \item \textbf{\mw{आप} \mw{मांस} \mw{खावो} \mw{हो}?} — \emph{āp mā̃s khāvo ho?} — Do you eat meat?
\end{itemize}

\begin{tcolorbox}[breakable,colback=teal!4,colframe=teal!35!black,title={Before you move on}]
What turns \textbf{\mw{वो} \mw{बीमार} \mw{है।}} into a question? (The voice, rising at the end.)
\end{tcolorbox}
\section[\texorpdfstring{\mw{हां} \mw{सा} · \mw{कोनी} \mw{सा} (hā̃ sā · konī sā)}{हां सा · कोनी सा (hā̃ sā · konī sā)}]{\mw{हां} \mw{सा}, \mw{कोनी} \mw{सा} — answering yes or no}

\label{lesson:MW-C158-answer}



\begin{quote}
Ask \emph{do you drink tea?} Then say the polite \emph{yes} and the plain \emph{no}: \textbf{\mw{हां} \mw{सा}}, \textbf{\mw{कोनी}}.
\end{quote}

\begin{grammarlens}[title={Grammar Lens: the answer, then the sentence}]
Answer with \textbf{\mw{हां} \mw{सा}} or \textbf{\mw{कोनी} \mw{सा}}, then say the sentence back:

\begin{quote}
\textbf{\mw{हां} \mw{सा}, \mw{हूं} \mw{पीऊं} \mw{हूं।}}
\end{quote}

\begin{quote}
\emph{hā̃ sā, hū̃ pīū̃ hū̃.} — Yes, I do.
\end{quote}

\begin{quote}
\textbf{\mw{कोनी} \mw{सा}, \mw{हूं} \mw{कोनी} \mw{पीऊं।}}
\end{quote}

\begin{quote}
\emph{konī sā, hū̃ konī pīū̃.} — No, I don't.
\end{quote}

\textbf{\mw{सा}} keeps the answer as polite as the question.
\end{grammarlens}

\subsection*{Your turn}
Say these aloud:

\begin{itemize}
\raggedright
  \item \textbf{\mw{राम} \mw{अठै} \mw{है}?} — \emph{rām aṭhai hai?} — Is Ram here?
  \item \textbf{\mw{कोनी} \mw{सा}, \mw{राम} \mw{अठै} \mw{कोनी।}} — \emph{konī sā, rām aṭhai konī.} — No, Ram isn't here.
  \item \textbf{\mw{हां} \mw{सा}, \mw{हूं} \mw{दूध} \mw{पीऊं} \mw{हूं।}} — \emph{hā̃ sā, hū̃ dūdh pīū̃ hū̃.} — Yes, I drink milk.
\end{itemize}

\begin{tcolorbox}[breakable,colback=teal!4,colframe=teal!35!black,title={Before you move on}]
The polite \emph{yes}? (\textbf{\mw{हां} \mw{सा}}.) The polite \emph{no}? (\textbf{\mw{कोनी} \mw{सा}}.)
\end{tcolorbox}
\section[\texorpdfstring{\mw{कद} (kad)}{कद (kad)}]{\mw{कद} (kad) — when}

\label{lesson:MW-C158-kad}



\begin{quote}
Before the new one: say \emph{where}, \emph{why} and \emph{who}. (\textbf{\mw{कठै}}, \textbf{\mw{क्यूं}}, \textbf{\mw{कुण}}.)
\end{quote}

\subsection*{You'll want to know: \mw{कद}}
\textbf{\mw{कद}} — \emph{kad} — \textquotedblleft{}when\textquotedblright{}.

\begin{quote}
\textbf{\mw{आप} \mw{कद} \mw{आवो} \mw{हो}?}
\end{quote}

\begin{quote}
\emph{āp kad āvo ho?} — When do you come?
\end{quote}

\begin{cousinweb}
\textbf{\mw{कद}} (\emph{kad}) asks \emph{when}. Monier-Williams gives Sanskrit \emph{kadā}, \textquotedblleft{}when?\textquotedblright{}, but CDIAL has no entry for \textbf{\mw{कद}}, so how the Marwari word came down is not settled. \textbf{\mw{कदी}}, \textquotedblleft{}sometimes, ever\textquotedblright{}, opens the same way.
\end{cousinweb}

\begin{grammarlens}[title={What you've built}]
\textbf{\mw{कद}} sits where the answer will go, as \textbf{\mw{कठै}} does: \textbf{\mw{वो} \mw{कद} \mw{जावै} \mw{है}?} (\emph{vo kad jāvai hai?}), when does he go?
\end{grammarlens}

\subsection*{Your turn}
\begin{itemize}
\raggedright
  \item \emph{Say it:} \emph{kad}
  \item \emph{Say it:} \textbf{\mw{आप} \mw{कद} \mw{आवो} \mw{हो}?} — \emph{āp kad āvo ho?} — When do you come?
  \item \emph{Recall:} say \emph{where}, then \emph{kad} again
\end{itemize}

\begin{tcolorbox}[breakable,colback=teal!4,colframe=teal!35!black,title={Before you move on}]
What does \textbf{\mw{कद}} mean? (\textquotedblleft{}When\textquotedblright{}.) Say it once more.
\end{tcolorbox}
\section[\texorpdfstring{\mw{ओळखणो} (oḷakhṇo)}{ओळखणो (oḷakhṇo)}]{\mw{ओळखणो} (oḷakhṇo) — to recognise}

\label{lesson:MW-C158-olakhno}



\begin{quote}
Before the new one: say \emph{to know} and \emph{to see}. (\textbf{\mw{जाणणो}}, \textbf{\mw{देखणो}}.) Then \emph{I understand a little}.
\end{quote}

\subsection*{You'll want to know: \mw{ओळखणो}}
\textbf{\mw{ओळखणो}} — \emph{oḷakhṇo} — \textquotedblleft{}to recognise\textquotedblright{}.

\begin{quote}
\textbf{\mw{हूं} \mw{यो} \mw{गीत} \mw{ओळखूं} \mw{हूं।}}
\end{quote}

\begin{quote}
\emph{hū̃ yo gīt oḷakhū̃ hū̃.} — I recognise this song.
\end{quote}

\begin{quote}
\textbf{\mw{आप} \mw{यो} \mw{गीत} \mw{ओळखो} \mw{हो}?}
\end{quote}

\begin{quote}
\emph{āp yo gīt oḷakho ho?} — Do you recognise this song?
\end{quote}

It opens on a letter you have not written yet: the vowel \emph{o}, standing alone. The next lesson puts it in your hand.

\begin{cousinweb}
\textbf{\mw{ओळखणो}} (\emph{oḷakhṇo}) goes back to Sanskrit \emph{upalakṣayati}, which CDIAL glosses \textquotedblleft{}beholds\textquotedblright{}. In that entry CDIAL gives Gujarati \emph{oḷakhvũ}, \textquotedblleft{}to know, distinguish\textquotedblright{}, and no Marwari form (CDIAL 2226).
\end{cousinweb}

\begin{grammarlens}[title={What you've built}]
Its stem \textbf{\mw{ओळख}} ends in a consonant: \textbf{\mw{ओळखूं} \mw{हूं}}, \textbf{\mw{ओळखै} \mw{है}}, \textbf{\mw{ओळखो} \mw{हो}}.
\end{grammarlens}

\subsection*{Your turn}
Say these aloud:

\begin{itemize}
\raggedright
  \item \emph{oḷakhṇo}
  \item \textbf{\mw{आप} \mw{यो} \mw{गीत} \mw{ओळखो} \mw{हो}?} — \emph{āp yo gīt oḷakho ho?} — Do you recognise this song?
  \item \textbf{\mw{हां} \mw{सा}, \mw{हूं} \mw{ओळखूं} \mw{हूं।}} — \emph{hā̃ sā, hū̃ oḷakhū̃ hū̃.} — Yes, I recognise it.
\end{itemize}

\begin{tcolorbox}[breakable,colback=teal!4,colframe=teal!35!black,title={Before you move on}]
What does \textbf{\mw{ओळखणो}} mean? (\textquotedblleft{}To recognise\textquotedblright{}.) Say it once more.
\end{tcolorbox}
\section[\texorpdfstring{\mw{ओ} (o)}{ओ (o)}]{\mw{ओ} (o) — the letter that opens \mw{ओळखणो}}

\label{lesson:MW-W48-o}



\begin{quote}
Before the new one: write \textbf{\mw{अ}} once, and name it. Then say \textbf{\mw{छोरो}} (\emph{chhoro}): its \emph{o} is the sign \textbf{\mw{ो}}.

Now say \textbf{\mw{ओळखणो}} (\emph{oḷakhṇo}), \textquotedblleft{}to recognise\textquotedblright{}. It opens on the vowel \emph{o} standing alone, and this lesson puts that letter in your hand.
\end{quote}

\subsection*{Script you'll notice: \mw{ओ}}
\textbf{\mw{ओ}} — \emph{o}.

It opens \textbf{\mw{ओळखणो}}. It is built on \textbf{\mw{अ}}: the joined bowls, the shoulder and the inner stem, then a trailing stem, as \textbf{\mw{आ}} has, and above it a separate arc, the same small arc the sign \textbf{\mw{ो}} carries.

\subsection*{Writing: \mw{ओ}}
\hlblockfigure{\detokenize{figures/MW-W48-o-filmstrip.pdf}}{How \mw{ओ} is written, stroke by stroke}

\begin{itemize}
\raggedright
  \item \textbf{1.} curve right around the upper bowl and continue down and around the lower bowl without lifting
  \item \textbf{2.} lift and sweep the middle shoulder right
  \item \textbf{3.} without lifting, climb up the inner stem to the headline
  \item \textbf{4.} descend the inner stem
  \item \textbf{5.} lift and descend the trailing stem
  \item \textbf{6.} lift at the trailing stem's headline junction and sweep the upper arc upward and left
  \item \textbf{7.} lift and draw the top shirorekhā left-to-right
\end{itemize}

\textbf{Pen lifts: 4.}

\begin{quote}
The source shows six movements; native writers draw \mw{ओ} as five pen-down strokes (60\% of HP Labs India's native-writer samples, 49 of 82). This ductus follows the six movements in order and direction, lifting four times: the middle shoulder runs on, climbing the inner stem and descending it, and every other movement starts with a lift.
\end{quote}

\begin{quote}
Stroke order is one attested teaching order; hands and teachers vary. Source: Saurmandal, ‘Devanagari \mw{ओ} stroke order.svg’, panels 1–6, Wikimedia Commons, 5 August 2023.
\end{quote}

\subsection*{Your turn}
\begin{itemize}
\raggedright
  \item \emph{Look:} at \textbf{\mw{ओळखणो}} and point at \mw{ओ} at its start
  \item \emph{Trace it:} \mw{ओ} three times, saying \emph{o} as you finish each one
  \item \emph{Write it:} \mw{अ}, then \mw{ओ}, from memory
  \item \emph{Read it:} \textbf{\mw{ओळखणो}} aloud once more
\end{itemize}

\begin{tcolorbox}[breakable,colback=teal!4,colframe=teal!35!black,title={Before you move on}]
Which letter is this — \mw{ओ}? (\textbf{o}.) Which word did it come from? (\textbf{\mw{ओळखणो}}, \emph{oḷakhṇo}, \textquotedblleft{}to recognise\textquotedblright{}.)
\end{tcolorbox}
\section[(dialogue)]{Yes or no? — a review}

\label{lesson:MW-R158-yes-or-no}



\begin{quote}
Ask \emph{do you drink tea?}, to \textbf{\mw{आप}}.
\end{quote}

\subsection*{Your turn}
\begin{itemize}
\raggedright
  \item \emph{Say it:} \emph{is Ram here?}; \emph{no, Ram isn't here}
  \item \emph{Say it:} \emph{do you eat meat?}; \emph{no, I don't eat meat}
  \item \emph{Say it:} \emph{when do you come?}; \emph{when does he go?}
  \item \emph{Say it:} \emph{do you recognise this song?}; \emph{yes, I recognise it}
  \item \emph{Read it:} \textbf{\mw{ओळखणो}}, and name its first letter
\end{itemize}

\begin{tcolorbox}[breakable,colback=teal!4,colframe=teal!35!black,title={Before you move on}]
\emph{When}? (\textbf{\mw{कद}}.) \emph{Yes}, politely? (\textbf{\mw{हां} \mw{सा}}.) \emph{No}, politely? (\textbf{\mw{कोनी} \mw{सा}}.)
\end{tcolorbox}
